
On
prend pour famille @@M0@@ la famille de tous les éléments de @@M1@@. Les hypothèses de 1.2 sont alors vérifiées, et par suite l'homomorphisme :
@@M2@@
est surjectif. D'après la définition de @@M3@@, ceci équivaut à dire que l'homomorphisme :
@@M4@@
est surjectif.

\textbf{Corollaire 1.4.} \uline{Soit} @@M5@@ \uline{un morphisme plat de schémas noethériens}. \uline{On suppose que pour tout point} @@M6@@ \uline{de} @@M7@@, \uline{l'homomorphisme canonique} @@M8@@ \uline{est un isomorphisme}. \uline{Alors l'homomorphisme} :
@@M9@@
\uline{est surjectif}. \uline{Si de plus} @@M10@@ \uline{possède une section} @@M11@@ \uline{de tor-dimension finie}, \uline{alors} @@M12@@ \uline{est un isomorphisme}.

On prend pour famille @@M13@@ la famille réduite au seul élément @@M14@@; on a alors @@M15@@, et, les conditions de 1.2 étant remplies, @@M16@@ est surjectif.

Si @@M17@@ admet une section @@M18@@ de tor-dimension finie, on peut définir un homomorphisme @@M19@@ (IV 2.12). Comme on a @@M20@@ d'après IV 2.12.1.1, @@M21@@ est injectif, donc est un isomorphisme.

\textbf{Remarque 1.5.} Pour que @@M22@@ soit de tor-dimension finie, il faut et il suffit que les fibres de @@M23@@ soient régulières aux points de @@M24@@, i.e.\ que @@M25@@ soit régulier aux points de @@M26@@ (ou encore lisse en ces points si on suppose @@M27@@
